\documentclass[11pt,a4paper]{article}
\usepackage[T1]{fontenc}
\usepackage[utf8]{inputenc}
\usepackage[french]{babel}
\usepackage{lmodern}
\usepackage{textcomp}
\usepackage{amsmath,amssymb}
\usepackage{geometry}
\usepackage{tabularx,array}
\usepackage[table]{xcolor}
\usepackage{tikz}
\geometry{left=8mm,right=8mm,top=7mm,bottom=7mm}
\setlength{\parindent}{0pt}
\pagestyle{empty}
\setlength{\tabcolsep}{3pt}
\renewcommand{\arraystretch}{1.05}
\renewcommand{\tabularxcolumn}[1]{m{#1}}

% --- Palette pensée pour une impression noir et blanc ---
\definecolor{canlight}{gray}{0.93}
\definecolor{canmid}{gray}{0.72}
\definecolor{candark}{gray}{0.22}
\definecolor{answerblue}{RGB}{0,82,165}

% --- Colonnes ---
\newcolumntype{Q}{>{\centering\arraybackslash}X}
\newcolumntype{N}{>{\centering\arraybackslash}m{6mm}}
\newcolumntype{R}{>{\centering\arraybackslash}m{22mm}}

% --- En-tête ---
\newcommand{\canheader}[1]{%
  \begin{center}
    {\LARGE\bfseries COURSE AUX NOMBRES}\par
    \vspace{0.7mm}
    {\color{canmid}\rule{27mm}{0.45pt}\hspace{2mm}\color{black}$\diamond$\hspace{1.4mm}{\large\bfseries #1}\hspace{1.4mm}$\diamond$\hspace{2mm}\color{canmid}\rule{27mm}{0.45pt}\color{black}}
  \end{center}
  \vspace{0.8mm}
  \renewcommand{\arraystretch}{1.25}%
  \begin{tabularx}{\linewidth}{>{\centering\arraybackslash}X >{\centering\arraybackslash}X >{\centering\arraybackslash}m{40mm}}
    \rowcolor{canlight}
    \footnotesize\textbf{30 questions} -- 8 minutes &
    \footnotesize\textbf{Sans calculatrice} &
    \footnotesize\textbf{Score :} \rule{13mm}{0.45pt}\,/30
  \end{tabularx}
  \renewcommand{\arraystretch}{1.05}%
  \vspace{1mm}
  \begin{center}
    {\scriptsize Réponse uniquement \quad$\diamond$\quad Reporter le score sur la fiche annuelle.}
  \end{center}
  \vspace{1.2mm}
}

\newcommand{\qrow}[2]{\textbf{#1} & \parbox[c][15.0mm][c]{\linewidth}{\centering #2} & \\ \hline}
\newcommand{\qrowfig}[2]{\textbf{#1} & \parbox[c][27mm][c]{\linewidth}{\centering #2} & \\ \hline}
\newcommand{\qrowcorr}[3]{\textbf{#1} & \parbox[c][15.0mm][c]{\linewidth}{\centering #2} & {\color{answerblue}\bfseries #3} \\ \hline}
\newcommand{\qrowfigcorr}[3]{\textbf{#1} & \parbox[c][27mm][c]{\linewidth}{\centering #2} & {\color{answerblue}\bfseries #3} \\ \hline}

\newcommand{\cantable}[1]{%
\arrayrulecolor{canmid}
\begin{tabularx}{\linewidth}{N|Q|R}
\rowcolor{canlight}\textbf{N\textsuperscript{o}}&\centering\textbf{Question}&\textbf{Réponse}\\\hline
#1
\end{tabularx}
\arrayrulecolor{black}
}

\newcommand{\hypotfigure}{%
\begin{tikzpicture}[baseline=(current bounding box.center),scale=0.72,line width=0.6pt]
  \node[font=\bfseries\normalsize] at (1.6,2.55) {Quelle est l'hypoténuse ?};
  \coordinate (A) at (0,0); \coordinate (B) at (0,1.85); \coordinate (C) at (3.2,0);
  \draw (A)--(B)--(C)--cycle;
  \draw (0,0) rectangle (0.32,0.32);
  \node[below left] at (A) {\small A};
  \node[above left] at (B) {\small B};
  \node[below right] at (C) {\small C};
\end{tikzpicture}}


\newcommand{\canheadercorr}[1]{%
  \begin{center}
    {\LARGE\bfseries COURSE AUX NOMBRES}\par
    \vspace{0.7mm}
    {\color{canmid}\rule{27mm}{0.45pt}\hspace{2mm}\color{black}$\diamond$\hspace{1.4mm}{\large\bfseries #1}\hspace{1.4mm}$\diamond$\hspace{2mm}\color{canmid}\rule{27mm}{0.45pt}\color{black}}
  \end{center}
  \vspace{0.8mm}
  \renewcommand{\arraystretch}{1.25}%
  \begin{tabularx}{\linewidth}{>{\centering\arraybackslash}X >{\centering\arraybackslash}X >{\centering\arraybackslash}m{40mm}}
    \rowcolor{canlight}
    \footnotesize\textbf{30 questions} -- 8 minutes &
    \footnotesize\textbf{Sans calculatrice} &
    \footnotesize\textcolor{answerblue}{\textbf{CORRIGÉ}}
  \end{tabularx}
  \renewcommand{\arraystretch}{1.05}%
  \vspace{1mm}
  \begin{center}
    {\scriptsize Les réponses apparaissent en \textcolor{answerblue}{\textbf{bleu}}.}
  \end{center}
  \vspace{1.2mm}
}


\begin{document}
\canheader{CAN 4}

\begin{minipage}[t]{0.493\textwidth}
\normalsize
\cantable{
\qrow{1}{$63\div7$}
\qrow{2}{$125+99$}
\qrow{3}{$(-6)-(-2)$}
\qrow{4}{Compléter : $3{,}7+\ldots=8$}
\qrow{5}{$1\,\%$ de 600}
\qrow{6}{La moitié du carré de 6}
\qrow{7}{Le plus grand : $2{,}04$ ou $2{,}4$ ?}
\qrow{8}{Donner sous forme de fraction le nombre qui, multiplié par 8, donne 3.}
\qrow{9}{$270{,}5\times0{,}01$}
\qrow{10}{Moyenne de 14 ; 9 ; 12 ; 5}
\qrow{11}{3 objets sur 12 sont rouges. Quelle fraction simplifiée des objets est rouge ?}
\qrow{12}{Quelle est la distance à zéro de $-3{,}5$ ?}
\qrow{13}{$\dfrac23$ de 27}
\qrow{14}{$3{,}25$ h = $\ldots$ h $\ldots$ min}
\qrow{15}{$54-4\times(18-8)$}
}
\end{minipage}\hfill
\begin{minipage}[t]{0.493\textwidth}
\normalsize
\cantable{
\qrow{16}{Simplifier $\dfrac{21}{35}$}
\qrow{17}{Aire d'un triangle de base 8 cm et de hauteur 5 cm}
\qrow{18}{Dans une série de 25 valeurs, 10 sont identiques. Leur fréquence décimale est :}
\qrow{19}{Écriture décimale de $\dfrac7{20}$}
\qrow{20}{Le périmètre d'une figure s'exprime en : cm, cm$^2$ ou cm$^3$ ?}
\qrow{21}{Parmi 705 ; 540 ; 1 340, lequel est divisible par 2, par 3 et par 5 ?}
\qrow{22}{$0{,}25\times28$}
\qrow{23}{45 L = $\ldots\,\mathrm{cm}^3$}
\qrow{24}{5 kg de melons coûtent 15 \texteuro. Combien coûtent 7 kg ?}
\qrow{25}{Un triangle rectangle a un angle aigu de $39^\circ$. L'autre angle aigu mesure :}
\qrow{26}{Parmi 50 ; 64 ; 72, quel nombre est un carré parfait ?}
\qrow{27}{Un rectangle de 3 cm sur 8 cm peut-il avoir un périmètre de 24 cm ?}
\qrow{28}{Le plus grand : $\dfrac56$ ou $\dfrac45$ ?}
\qrow{29}{Dans $3x$, le symbole caché entre 3 et $x$ est : $+$ ou $\times$ ?}
\qrow{30}{Pour $x=4$, l'égalité $2x+1=9$ est-elle vraie ?}
}
\end{minipage}

\clearpage

\canheadercorr{CAN 4}

\begin{minipage}[t]{0.493\textwidth}
\normalsize
\cantable{
\qrowcorr{1}{$63\div7$}{$9$}
\qrowcorr{2}{$125+99$}{$224$}
\qrowcorr{3}{$(-6)-(-2)$}{$-4$}
\qrowcorr{4}{Compléter : $3{,}7+\ldots=8$}{$4{,}3$}
\qrowcorr{5}{$1\,\%$ de 600}{$6$}
\qrowcorr{6}{La moitié du carré de 6}{$18$}
\qrowcorr{7}{Le plus grand : $2{,}04$ ou $2{,}4$ ?}{$2{,}4$}
\qrowcorr{8}{Donner sous forme de fraction le nombre qui, multiplié par 8, donne 3.}{$\dfrac38$}
\qrowcorr{9}{$270{,}5\times0{,}01$}{$2{,}705$}
\qrowcorr{10}{Moyenne de 14 ; 9 ; 12 ; 5}{$10$}
\qrowcorr{11}{3 objets sur 12 sont rouges. Quelle fraction simplifiée des objets est rouge ?}{$\dfrac14$}
\qrowcorr{12}{Quelle est la distance à zéro de $-3{,}5$ ?}{$3{,}5$}
\qrowcorr{13}{$\dfrac23$ de 27}{$18$}
\qrowcorr{14}{$3{,}25$ h = $\ldots$ h $\ldots$ min}{3 h 15 min}
\qrowcorr{15}{$54-4\times(18-8)$}{$14$}
}
\end{minipage}\hfill
\begin{minipage}[t]{0.493\textwidth}
\normalsize
\cantable{
\qrowcorr{16}{Simplifier $\dfrac{21}{35}$}{$\dfrac35$}
\qrowcorr{17}{Aire d'un triangle de base 8 cm et de hauteur 5 cm}{20 cm$^2$}
\qrowcorr{18}{Dans une série de 25 valeurs, 10 sont identiques. Leur fréquence décimale est :}{$0{,}4$}
\qrowcorr{19}{Écriture décimale de $\dfrac7{20}$}{$0{,}35$}
\qrowcorr{20}{Le périmètre d'une figure s'exprime en : cm, cm$^2$ ou cm$^3$ ?}{cm}
\qrowcorr{21}{Parmi 705 ; 540 ; 1 340, lequel est divisible par 2, par 3 et par 5 ?}{$540$}
\qrowcorr{22}{$0{,}25\times28$}{$7$}
\qrowcorr{23}{45 L = $\ldots\,\mathrm{cm}^3$}{$45\,000$ cm$^3$}
\qrowcorr{24}{5 kg de melons coûtent 15 \texteuro. Combien coûtent 7 kg ?}{21 \texteuro}
\qrowcorr{25}{Un triangle rectangle a un angle aigu de $39^\circ$. L'autre angle aigu mesure :}{$51^\circ$}
\qrowcorr{26}{Parmi 50 ; 64 ; 72, quel nombre est un carré parfait ?}{$64$}
\qrowcorr{27}{Un rectangle de 3 cm sur 8 cm peut-il avoir un périmètre de 24 cm ?}{non (22 cm)}
\qrowcorr{28}{Le plus grand : $\dfrac56$ ou $\dfrac45$ ?}{$\dfrac56$}
\qrowcorr{29}{Dans $3x$, le symbole caché entre 3 et $x$ est : $+$ ou $\times$ ?}{$\times$}
\qrowcorr{30}{Pour $x=4$, l'égalité $2x+1=9$ est-elle vraie ?}{oui}
}
\end{minipage}
\end{document}
